\section*{Project contract}
Implement \texttt{eval\_expr} for a byte slice containing an integer expression.
Pass its full address in \texttt{a0} and its length, zero through 64 bytes, in
\texttt{a1}. The slice lies within one of the supplied 96-byte input arenas and
may begin at an unaligned byte address. Its length is authoritative: no trailing
zero byte is required, and a zero byte inside the slice is an invalid input byte.
The function produces no console output and writes no caller data, even temporarily.
The selected slice is the only caller-data input: do not read the other arena, metadata or guards. Submission-owned constant tables may be read. Private stack storage is allowed. Preserve all bits of \texttt{s0} through
\texttt{s7}, \texttt{fp}, and \texttt{sp}, and return through the supplied
\texttt{ra}. Source labels and temporary-register choices inside the function are
not part of the free-writing interface.

Return a status in \texttt{v0}. Zero means success, with the signed 32-bit value
in \texttt{v1}. Minus one means syntax error, minus two means arithmetic overflow,
and minus three means division by zero. On error, \texttt{v1} instead contains
a zero-based byte offset. An offset equal to the supplied length identifies a
missing item at end of input. Every input load must fit wholly inside the supplied
slice. Reading ahead within that slice is permitted when it does not change the
specified result or diagnostic.

The language accepts decimal digit sequences, binary addition, subtraction,
multiplication and division, unary plus and minus, and parentheses. Space, horizontal
tab, carriage return and line feed may appear between tokens. They do not join
separated digit sequences into one number. No comments, identifiers, implicit
multiplication, hexadecimal literals or floating-point notation are accepted.
Every decimal literal has a value from zero through 2,147,483,647. Each intermediate
arithmetic result must lie from $-2^{31}$ through $2^{31}-1$. Division truncates its
quotient toward zero. The order of validation and evaluation below determines which
failure is reported when an input permits more than one diagnosis.
